% Part: sets-functions-relations
% Chapter: sets
% Section: important-sets

\documentclass[../../../include/open-logic-section]{subfiles}

\begin{document}

\olfileid{sfr}{set}{imp}
\olsection{一些重要的集合}

\begin{ex}
我们将主要处理元素为数学对象的集合。其中有四个集合足够重要，因而有专门的名称：
\begin{multline*}
    \Nat = \{0, 1, 2, 3, \ldots\} \\
    \shoveright{\text{自然数集}}\\
    \shoveleft{\Int = \{\ldots, -2, -1, 0, 1, 2, \ldots\}} \\
    \shoveright{\text{整数集}}\\
    \shoveleft{\Rat = \Setabs{\nicefrac{m}{n}}{m, n \in \Int\text{、}n \neq 0}}\\
    \shoveright{\text{有理数集}}\\
    \shoveleft{\Real = (-\infty, \infty)}\\
    \text{以及实数集（连续统）}
\end{multline*}。
这些都是\emph{无限}集，也就是说，它们每个都有无穷多个元素。

顺着这些集合的序列，我们实际上是在将\emph{更多}的数加入我们的储备。
事实上，显然有 $\Nat \subseteq \Int \subseteq \Rat \subseteq \Real$：毕竟，每个自然数都是整数，每个整数都是有理数，每个有理数都是实数。
同样显然的是 $\Nat \subsetneq \Int \subsetneq \Rat$，因为 $-1$ 是整数而不是自然数，且 $\nicefrac{1}{2}$ 是有理数而不是整数。
不太明显的是 $\Rat \subsetneq \Real$，即存在某些实数不是有理数\oliflabeldef{sfr:arith:real:realline}{，但我们将在\olref[arith][real]{realline}中回到这个问题}{}。

我们有时也会用到正整数集 $\PosInt = \{1,
2, 3, \dots\}$ 以及仅由前两个自然数组成的集合 $\Bin = \{0, 1\}$。
\end{ex}

\begin{tagblock}{compsci}
\begin{ex}[字符串]
另一个有趣的例子是字母表 $A$ 上所有\emph{有限字符串}的集合 $A^{*}$：任何由 $A$ 中元素组成的有限序列都是 $A$ 上的字符串。对于每个字母表~$A$，$A$ 上的字符串都包含\emph{空串 $\Lambda$}。例如，
\begin{multline*}
\Bin^*
=\{\Lambda,0,1,00,01,10,11,\\
000,001,010,011,100,101,110,111,0000,\ldots\}.
\end{multline*}
如果 $x=x_{1}\ldots x_{n}\in A^{*}$是由 $A$ 中 $n$ 个“字母”组成的字符串，则我们说该字符串的\emph{长度}为~$n$，记作 $\len{x}=n$。
\end{ex}
\end{tagblock}

\begin{ex}[无限序列]
对于任意集合 $A$，我们也可以考虑由 $A$ 中元素组成的无限序列的集合~$A^\omega$。无限序列 $a_1a_2a_3a_4\dots$ 由一个单向无限的对象列表组成，其中每个对象都是 $A$ 的一个元素。
\end{ex}

\end{document}